\chapter{Introduction}
\label{ch:intro}

Charged particles moving in electromagnetic fields do not travel in straight lines. Depending on the field geometry, a particle can gyrate around a field line, drift slowly across it or bounce back and forth between different regions of the field. Understanding this motion is fundamental to space plasma physics, with direct relevance to energetic particles trapped in planetary radiation belts such as Earth's Van Allen belts. Particle orbit theory provides the analytic framework for describing these dynamics in idealised geometries \citep{boyd2003}.

These analytic results rely on symmetry. Once the field becomes spatially non-uniform, curved or time-dependent, closed-form trajectories are generally no longer available and the Lorentz force equation must be solved numerically. Numerical integration introduces truncation errors that accumulate over time, so any solver must be carefully verified before its output can be trusted.

This project builds and validates this framework. The equations of motion are solved with an adaptive Runge--Kutta scheme
(\texttt{scipy.integrate.solve\_ivp}, RK45), with different field configurations implemented as interchangeable modules so the same integrator handles every test without modification. The solver is verified through a range of test cases against known analytic results: circular gyromotion, the $\mathbf{E} \times \mathbf{B}$,
gradient-$B$ and curvature drifts, and magnetic mirror effects. In every case, the numerical output is compared quantitatively against the corresponding analytic formula and kinetic energy conservation is monitored.

The validated solver is then applied to dipole field configurations to simulate particle dynamics in planetary magnetospheres \citep{baumjohann1996}. Full orbit simulations demonstrate bounce motion between mirror points and confirm approximate conservation of the magnetic moment $\mu$. Numerically measured bounce periods are compared against analytic values from the standard bounce-period integral. A
separate guiding-centre integrator is run alongside the full Lorentz solver to test directly when the adiabatic approximation holds. The framework is also applied in physical SI units for a $10\,$keV electron at $L = 3$, representative of Earth's radiation belts, and extended to tilted dipole fields representative of Neptune ($47^\circ$) and Uranus ($59^\circ$), where the displaced magnetic equator produces asymmetric bounce motion in geographic coordinates.

The report is aimed at an Honours mathematics or physics student familiar with classical mechanics and vector calculus.
